% 第8章 8.3.3：常值族 f_a=a，a∈[0,M]，M=1，精度 eps=0.15，一致距离。
% 每个代表 c 只覆盖长度 2eps=0.3 的参数区间 [c-eps,c+eps]，故至少需要 ceil(M/(2eps))=4 个代表；
% 均匀中心 (2j-1)M/(2N)=0.125,0.375,0.625,0.875（N=4）达到这一数量。示意图，坐标按比例。
\begin{tikzpicture}[font=\figurefont\footnotesize,x=80mm,y=1mm]
\draw[FigureStroke,line width=.7pt] (-0.05,0)--(1.08,0);
\fill[FigureBlueFill] (0,-1.2) rectangle (1,1.2);
\draw[FigureInk,line width=1.2pt] (0,0)--(1,0);
\foreach \x/\l in {0/0,1/M}{\draw[FigureInk,line width=.6pt] (\x,-1.6)--(\x,1.6); \node[anchor=north,font=\figurefont\scriptsize] at (\x,-2) {$\l$};}
\foreach \c/\h [count=\j] in {0.125/4,0.375/8,0.625/4,0.875/8}{
  \draw[FigureBlue,line width=1.6pt] (\c-0.15,\h)--(\c+0.15,\h);
  \draw[FigureBlue,line width=.5pt] (\c-0.15,\h-1)--(\c-0.15,\h+1); \draw[FigureBlue,line width=.5pt] (\c+0.15,\h-1)--(\c+0.15,\h+1);
  \fill[FigureGreen] (\c,0) circle (1.4pt);
  \draw[FigureMuted,line width=.4pt,dashed] (\c,0)--(\c,\h);
  \node[anchor=north,font=\figurefont\scriptsize] at (\c,-2) {$c_{\j}$};}
\draw[<->,FigureInk,line width=.5pt] (0.725,11)--node[above,font=\figurefont\scriptsize]{$2\varepsilon$}(1.025,11);
\node[anchor=west,font=\figurefont\scriptsize] at (1.09,0) {参数 $a$};
\end{tikzpicture}
